\documentclass[11pt,a4paper]{article}
\usepackage[T1]{fontenc}
\usepackage{lmodern}
\usepackage{isabelle,isabellesym}
\usepackage{amsmath,amssymb}
\usepackage{pdfsetup}

\isabellestyle{it}
\renewcommand{\isastyletext}{\normalsize\normalfont}
\renewcommand{\isastyletxt}{\normalfont}

\newcommand{\ket}[1]{|#1\rangle}
\newcommand{\norm}[1]{\|#1\|}
\newcommand{\C}{\mathbb{C}}
\newcommand{\R}{\mathbb{R}}
\DeclareMathOperator{\diag}{diag}

\begin{document}

\title{The Functional Calculus of Hermitian Matrices and Product Formulas}
\author{Dustin Bryant}
\maketitle

\begin{abstract}
For a Hermitian matrix $A$ and a function $g:\R\to\C$ we define $g(A)$ as the matrix
that multiplies every eigenvector of $A$ with eigenvalue $\lambda$ by $g(\lambda)$, and
derive the algebra of this functional calculus: $g(A)h(A)=(gh)(A)$, $g(A)^*=\bar g(A)$,
unitarity for $|g|=1$ on the spectrum, the time evolution $e^{itA}$ and the inverse.
On top of it we verify the Lie--Trotter product formula
$\norm{(e^{itH_1/r}\cdots e^{itH_m/r})^r-e^{it(H_1+\dots+H_m)}}\le 2(mLt)^2/r$ and
Suzuki's higher-order formulas: Suzuki's recursion $S_{2k+2}$ agrees with
$e^{itH}$ up to degree $2k+2$ of the Taylor expansion, and $r$ steps of it are within
$r(2|F|+1)(\Lambda|t|/r)^{2k+3}e^{\Lambda|t|/r}/(2k+3)!$ of $e^{itH}$.  The matrices are
those of the AFP entry \emph{Jordan Normal Form} \cite{thiemann2015jnf}; the spectral
theorem comes from the Schur decomposition of the AFP entry \emph{Quantum Hoare Logic}
\cite{liu2019afp}, the adjoint, unitarity and the vector norm from
\emph{Isabelle Marries Dirac} \cite{bordg2020afp}.  The product formulas are due to
Trotter \cite{trotter1959} and Suzuki \cite{suzuki1991}; they are the basis of quantum
simulation \cite{lloyd1996,berry2007}.
\end{abstract}

\tableofcontents
\newpage

\parindent 0pt\parskip 0.5ex

\input{session}

\bibliographystyle{abbrv}
\bibliography{root}

\end{document}
